\section{Threats to Validity}
\label{sec:threats}

The threats taxonomy follows Wohlin et al.~\cite{wohlin2012experimentation}.

\textbf{Conclusion validity.} The six experiments comprise 22{,}500
P@50 observations summarised with 95\% percentile bootstrap CIs
(10{,}000 resamples). Some hypothesis decisions depend on small
sample Spearman correlations (E4 H1/H2/H4 on 5 $\tau$ points;
E3/E5 H3 on 3 K points); we report point values and treat them as
descriptive rather than strongly inferential. The E6 H2 fine-margin
rejection ($+0.0090$ vs required $+0.010$) is within the bootstrap
CI for the K=200 cell-mean comparison; we report the binary
rejection per protocol and the substantive direction as the
contribution.

\textbf{Internal validity.}

\textit{(i) Selection-policy coupling.} All six experiments share
the convention that HygienePrio-full under fixed weights drives the
fleet trajectory. The recalibration decision affects scoring
weights only, not fleet evolution. In a closed-loop deployment the
selection policy would also be recalibrated; the experiments would
not generalise directly. Paper~9~\cite{paper9selftraj} addressed the
cross-method trajectory question separately and showed that
relaxing the fixed-trajectory convention re-attributes Paper~6's
K=200 absolute collapse to selection coupling. The present paper
inherits the fixed-trajectory convention and bounds its claims to
that operating regime.

\textit{(ii) Pre-registered parameters not tuned.} Procedure
parameters (E2: EWMA $\alpha = 0.6$, history $H = 3$; E3: $\tau =
0.05$; E5: $\tau_K = \{0.20, 0.05, 0.02\}$; E6: $k = 0.04, h =
0.10$) were locked before evaluation-seed inspection. Other
parameter choices could shift the H1/H2 trade-off curve; sensitivity
sweeps over these parameters are future work for each detector
family.

\textit{(iii) Selection of detector families.} The six experiments
sample the simplest-deployable end of the detector design space.
Two-sided CUSUM, EWMA-CUSUM, Bayesian online change-point methods,
and capacity-aware CUSUM are out of scope; their results could
shift any of the per-experiment conclusions.

\textbf{Construct validity.} P@50 at $K = 200$ becomes noisy at
late windows due to small positive-set sizes; the cell-mean
comparisons partially mitigate this. CUSUM's accumulator state
introduces path-dependence in the strategy's behaviour that
single-window detectors do not have; the E6 K=200 gain depends on
the specific $|\Delta_w|$ time series, including the W4 noise
excursion that the synthetic generator produces. Real fleet
calibration-target time series may have different temporal
structure; the qualitative claim (CUSUM filters single-window noise)
is more robust than the specific $+0.9$~pp gain magnitude.

\textbf{External validity.} Inherited synthetic EEHDA priors,
fixed $\lambda$, 14-day window, 0.15 patch-debt reduction, $\sim$8\%
KEV prevalence. Real fleets with different rates of per-window
distributional shift would re-locate the H1/H2 trade-off curves and
the CUSUM gain magnitude. No claim is made about generalisation to
real enterprise fleets in the absence of external validation on
real exploitation-event data.
